\section{Adversary Model and Claim Boundary}
\label{sec:extended-adversary}

\subsection{System instance and protected transition}

We model an attachment as an untrusted tuple
$a=(x,n,m,d)$, where $x\in\{0,\ldots,255\}^{*}$ is a finite byte string,
$n$ is a caller-supplied display name, $m$ is an optional supplied media type,
and $d\in\{\mathsf{outgoing},\mathsf{incoming}\}$ identifies the endpoint
direction. Outgoing processing occurs before encryption. Incoming processing
occurs after authenticated decryption and before an application parser or
renderer. In both cases the protected transition is from attacker-influenced
plaintext to a locally publishable media representation.

The encrypted relay is outside the plaintext decision boundary. It may delay,
drop, replay, reorder, or replace ciphertext subject to the guarantees of the
messaging protocol, but it need not receive attachment plaintext and is not
trusted to authorize publication. This placement preserves the intended
end-to-end confidentiality architecture. It does not extend the cryptographic
security definition of the messaging protocol. The contribution considered
here begins when plaintext bytes are available at an authenticated endpoint and
ends when a reconstructed artifact is either published or withheld.

Let $p\in\mathcal{P}$ denote a versioned local policy and let
$\mathcal{K}_{p}$ be its finite set of admitted format profiles. A policy binds
name rules, media-type aliases, routing rules, canonical output languages,
reconstruction requirements, signature rules, and resource limits. The
authenticated application binary, the selected policy, and the canonical
profile registry are trusted inputs to the decision. Attachment bytes and all
descriptive fields are not trusted. This distinction follows the general
language-theoretic requirement that recognizers define the language accepted at
a trust boundary rather than infer safety from informal labels alone
\cite{momot2016}.

\subsection{Adversary capabilities}

The adversary is adaptive and implementation-aware. It may know the source
code, dependency versions, canonical grammars, resource limits, and the active
policy. It may choose each byte of $x$, including length fields,
offsets, checksums, compressed payloads, dimensions, sample tables, frame
counts, timestamps, metadata, and terminal markers. It may choose misleading or
malformed values for $n$ and $m$. It may submit many related inputs and use
observable outcomes to refine a later submission. Security therefore cannot
depend on format obscurity, unpredictable parser selection, or a secret
signature list.

The attacker may construct prefix camouflage, appended data, overlapping
grammars, and nested active objects. It may place a valid media signature after
an executable prefix, append a second grammar after a valid media terminator,
or construct bytes accepted differently by two parsers. It may also place an
attachment, script, foreign track, or executable object in a location that a
permissive container grammar would ignore. Because arbitrary substrings can
occur coincidentally in compressed media, the model does not equate a byte
substring with a parsed embedded object. A nested-object claim requires a
location-aware parser or an exact indicator whose scope is fixed by policy.

Resource exhaustion is part of the core model. The adversary may declare very
large dimensions, durations, nesting depths, component counts, or output sizes.
It may provide inputs with high decompression ratios, malformed arithmetic, or
codec behavior that approaches a deadline. It may also trigger decoder errors,
child-process launch failures, timeouts, cancellations, and panics in reachable
code. Repeated requests may amplify any per-request cost. The method therefore
treats availability limits as acceptance conditions rather than operational
advice.

\begin{table}[t]
\centering
\caption{Adversary capability and the corresponding modeled defense. A row
identifies a measurable carrier property, not a behavioral malware-family
classifier.}
\label{tab:adversary-capabilities}
\small
\begin{tabularx}{\textwidth}{L{0.24\textwidth}YY}
\toprule
Attacker control & Security-relevant effect & Modeled response \\
\midrule
Name and supplied media type & Misrouting, hidden executable suffix, or type
confusion & Portable-name grammar and evidence convergence before dispatch \\
Arbitrary byte layout & Prefix, suffix, malformed boundary, or overlapping
grammar & Offset-zero parsing, exact terminal offset, and ambiguity rejection \\
Container elements and metadata & Nested active object, stale offset, or
privacy-bearing field & Location-aware component allowlist and reconstruction \\
Compressed media payload & Decoder trigger or retained covert information &
Semantic reconstruction when required, otherwise a bounded structural claim \\
Counts, dimensions, and ratios & Excessive allocation, expansion, or execution
time & Pre-operation limits, cancellation, deadline, and fail-closed errors \\
Repeated submissions & Availability pressure and adaptive probing & Per-request
bounds, deterministic policy, and auditable terminal outcomes \\
\bottomrule
\end{tabularx}
\end{table}

\subsection{Trust assumptions}

The analytical arguments assume that the endpoint begins processing in an
uncompromised state. The policy and alias registry are authenticated and loaded
without attacker modification. The implementation executes the intended code,
cryptographic digests are collision resistant for provenance purposes, and the
operating system enforces process and file-system primitives according to their
documented semantics. When an external decoder is invoked, its output is still
untrusted until post-reconstruction validation succeeds. Semantic
reconstruction nonetheless depends on the decoder producing the semantic state
that the encoder consumes. Process supervision can bound some effects of decoder failure,
but invoking a decoder does not make that decoder correct.

Resource enforcement must occur before the operation it bounds. Arithmetic for
sizes and offsets must reject overflow before conversion. Wall-clock limits
depend on an operating-system mechanism capable of terminating or abandoning
the bounded operation. These obligations are deferred to the controlled fault
and fuzz campaigns.

For filesystem publication, the destination directory and its namespace are
trusted against concurrent hostile mutation during one transaction. The file
system is assumed to provide same-directory atomic namespace publication for
the no-clobber persistence primitive. The current implementation synchronizes
the temporary file before persistence. It does not synchronize the parent
directory after the namespace operation. The formal publication claim is
therefore atomic visibility of a newly published path, not guaranteed survival
of that directory entry across every power-loss model. This boundary is made
explicit in \cref{sec:publication-semantics}.

\subsection{Threat strata and non-goals}

Viruses, worms, Trojans, ransomware, spyware, downloaders, droppers, and stagers
motivate threat scenarios, but their behavioral labels are not classifier
outputs. A file-infector virus can motivate an executable prefix or appended
host modification. A downloader retrieves another payload after execution. A
dropper extracts or installs a payload already carried in the initial object. A
stager prepares a later execution phase without requiring either mechanism.
A worm scenario requires self-propagation after execution and can add repeated
delivery context. Spyware is used only when independent evidence labels a
behavioral payload. Passive location, author, device, or timestamp metadata is
placed in the privacy-leakage stratum and is not relabeled as spyware. In every
case the measured system outcome concerns a carrier property such as grammar
disagreement, forbidden structure, metadata presence, or reconstruction level.
It does not establish that the system recognized runtime behavior. This usage is
consistent with using conventional malware categories for incident analysis
while preserving a narrower file-processing claim \cite{nist80083}.

Several goals are outside the model. \AFD{} does not prove that arbitrary
accepted bytes are harmless. It does not identify phishing meaning, abusive
imagery, deceptive text rendered into pixels, or malicious information encoded
in legitimate samples. Structural reconstruction does not neutralize content
inside a retained pixel, audio, or codec bitstream. Even semantic reconstruction
does not prove the decoded meaning benign. It changes the carrier by deriving a
new encoding from bounded semantic state. The method does not replace endpoint
sandboxing, operating-system exploit mitigations, antimalware engines, or
application-specific authorization.

Compromise before processing, modification of the authenticated binary or
policy, and malicious behavior in a privileged local component are outside the
claim. Side channels beyond the terminal outcome and coarse resource behavior
are not analyzed. The following propositions are analytical arguments, not
machine-checked proofs or substitutes for deferred empirical evidence.
